% 原PDF物理页155–172：第六章，完整翻译。
\SourcePage{155}
本章旨在说明并论证一些方法，用它们可以从有限域的代数闭包上的命题（例如5.3和5.4中的命题），推出$\C$上的命题。

\section*{6.1. 原理}
\addcontentsline{toc}{section}{6.1. 原理}
\par\medskip\noindent\textbf{6.1.1.} 要从$\overline{\F}$上的命题推出关于复代数簇的几何命题，需要分两步进行；中间环节是一个$\C$上的命题，但这个命题不涉及基域$\C$的拓扑。如果所讨论的是\term{cohomology}{上同调}或\term{homotopy}{同伦}命题，就把复代数簇的经典拓扑换成\term{etale}{平展}拓扑，并把系数$\Z/\ell,\Z,\Q,\C$分别换成$\Z/\ell,\Z_\ell,\Q_\ell,\overline{\Q}_\ell$，从而得到这个中间命题。这种转换并非总能进行：
\begin{itemize}
\item \term{etale}{平展}拓扑探测不到$\pi_1$本身，而只能探测其预有限完备化$\widehat{\pi}_1$（至少对正规簇如此）。若$\pi_1$有限，它还能探测$\widehat{\pi}_i=\pi_i\otimes\widehat{\Z}$（$i>1$）；但若$\pi_1$无限，高阶$\pi_i$（它们未必有限生成）的信息就可能丢失，参见Artin--Mazur，\emph{Etale homotopy}，Springer Lecture Notes in Math. 100。
\item \term{etale}{平展}拓扑只有对\term{l-adic}{$\ell$-进}系数才具有良好性质。不能取$\Z,\C$或$\R$为系数：整格和二次型的符号差等信息因此丢失。
\item \term{cohomology}{上同调}的Hodge结构（或\term{mixed}{混合}Hodge结构）在\term{l-adic}{$\ell$-进}\term{cohomology}{上同调}中丢失。不过，$H^*(X,\Q)$的\term{weight-filtration}{权滤过}有几何定义，因此有相应的\term{l-adic}{$\ell$-进}版本。应用见6.2.3。
\end{itemize}
此外，所讨论的簇和映射必须是复\emph{代数}的，只有复解析性还不够。例如，5.3和5.4的结果将给出射影簇上的一些定理，而它们在Kähler簇上的类似命题尚不清楚。

\par\medskip\noindent\textbf{6.1.2.} 设$X$为$\C$上的\term{finite-type}{有限型}\term{scheme}{概形}。连接拓扑空间$X(\C)$与$X$的\term{etale}{平展}拓扑，有一个比较态射（即一个\term{topos}{拓扑斯}态射）
\SourcePage{156}
\[
\varepsilon:X(\C)\longrightarrow X_{\mathrm{\acute et}}.
\]
基本结果如下。

\noindent\textbf{(A)} 函子$\varepsilon^*$给出$X_{\mathrm{\acute et}}$上的\term{constructible}{可构造}集合\term{sheaf}{层}范畴与$X(\C)$上的\term{constructible}{可构造}集合\term{sheaf}{层}范畴之间的\term{equivalence}{等价}。这里“\term{constructible}{可构造}”指在某个\emph{代数}\term{stratification}{分层}的每个\term{stratum}{层片}上局部常值且纤维有限。还有以下变体：

\noindent\textbf{(A$'$)} 对有限环$R$，有范畴\term{equivalence}{等价}
\[
\varepsilon^*:\{X_{\mathrm{\acute et}}\text{上的\term{constructible}{可构造}$R$-模\term{sheaf}{层}}\}
\xrightarrow{\sim}\{X(\C)\text{上的\term{constructible}{可构造}$R$-模\term{sheaf}{层}}\}.
\]
\noindent\textbf{(A$''$)} 对任意素数$\ell$，有范畴\term{equivalence}{等价}
\[
\varepsilon^*:\{X_{\mathrm{\acute et}}\text{上的\term{constructible}{可构造}$\Z_\ell$-\term{sheaf}{层}}\}
\xrightarrow{\sim}\{X(\C)\text{上的\term{constructible}{可构造}$\Z_\ell$-模\term{sheaf}{层}}\}.
\]
对于$\Q_\ell$-\term{sheaf}{层}，$\varepsilon^*$只有\term{fully-faithful}{全忠实}性。设$X$为正规连通簇，$x\in X(\C)$为基点，$F$为有限秩的局部常值$\Q_\ell$-向量空间\term{sheaf}{层}。那么，$F$属于其\term{essential-image}{本质像}，当且仅当$\pi_1(X(\C),x)$在$F_x$上的作用保持某个格，即保持$F_x$中生成$F_x$的某个有限生成$\Z_\ell$-子模。若$F$在某个\term{stratification}{分层}$\mathcal S$（按2.1的含义）的每个正规连通\term{stratum}{层片}上局部常值，则要求每个$F|S$（$S\in\mathcal S$）都满足同一条件。$E_\lambda$-\term{sheaf}{层}和$\overline{\Q}_\ell$-\term{sheaf}{层}的情形亦然。不过，由于这些$\pi_1$都是有限生成群，若$F$为$\Q$-向量空间的\term{constructible}{可构造}\term{sheaf}{层}，则对几乎所有$\ell$，$F\otimes\Q_\ell$都属于$\varepsilon^*$的\term{essential-image}{本质像}。研究\term{constructible}{可构造}\term{sheaf}{层}时，与其借助这一观察归约到\term{l-adic}{$\ell$-进}情形，往往不如把待证命题归约到(A$'$)更方便。例子见6.2.1。

\term{derived-category}{导出范畴}也有类似命题：

\noindent\textbf{(B$'$)} 对有限环$R$，有范畴\term{equivalence}{等价}
\[
\varepsilon^*:D_c^b(X,R)\xrightarrow{\sim}D_c^b(X(\C),R).
\]
下标$c$表示“\term{cohomology}{上同调}\term{sheaf}{层}\term{constructible}{可构造}”。

\noindent\textbf{(B$''$)} 有
\[
\varepsilon^*:D_c^b(X,\Z_\ell)\xrightarrow{\sim}D_c^b(X(\C),\Z_\ell).
\]
左边是定义为$2\text{-}\varprojlim D_{ctf}^b(X,\Z/\ell^n)$的\term{triangulated-category}{三角范畴}（t.f.表示有限Tor维数）；右边则是通常\term{derived-category}{导出范畴}的一个子范畴。
\SourcePage{157}
对于$\Q_\ell$-\term{cohomology}{上同调}，也仍然只有\term{fully-faithful}{全忠实}性的结果。$K$属于\term{essential-image}{本质像}，当且仅当它的各个\term{cohomology}{上同调}\term{sheaf}{层}都属于\term{essential-image}{本质像}。$E_\lambda$-\term{cohomology}{上同调}与$\overline{\Q}_\ell$-\term{cohomology}{上同调}的情形亦然。

函子$\varepsilon^*$与通常的\term{cohomology}{上同调}运算交换，至少在限于\term{constructible}{可构造}\term{sheaf}{层}时如此。关键结果是SGA~4 XVI 4.1：

\noindent\textbf{(C)} 设$f:X\to Y$为$\C$上\term{finite-type}{有限型}\term{scheme}{概形}之间的态射。对$X$上任意\term{constructible}{可构造}\term{torsion}{挠}\term{sheaf}{层}$F$，有
\[
\varepsilon^*R^qf_*F\xrightarrow{\sim}R^qf_*(\varepsilon^*F).
\]
这里$F$的\term{constructible}{可构造}性不可缺少，如下例所示：$X$为\term{affine}{仿射}直线，$Y$为一点，$F=\bigoplus_{n\in\Z}(\Z/\ell\Z)_n$。

由此形式地推出：

\noindent\textbf{(C$'$)} 在$\Z/\ell^n,\Z_\ell,\Q_\ell$或$\overline{\Q}_\ell$-\term{cohomology}{上同调}中，也就是说，在相应的有界\term{constructible}{可构造}\term{derived-category}{导出范畴}中，$\varepsilon^*$与函子$Rf_*,Rf_!,f^*,Rf^!$以及$\overset L\otimes$和$R\mathcal Hom$交换。

\par\medskip\noindent\textbf{证明。} 关键结果是SGA~4 XVI 4.1，即$R^qf_*$的比较定理。对\term{constructible}{可构造}集合\term{sheaf}{层}（分别地，群\term{sheaf}{层}），当$q=0$（分别地，$q=0$或$1$）时也成立。对局部常值集合\term{sheaf}{层}，(A$'$)来自$\pi_0$和$\pi_1$的比较定理；后者由SGA~4 XVI 4.1取$q=0,1$和$S=\Spec(\C)$得到。一般情形用\term{noetherian}{诺特}归纳法处理：取足够小的稠密开集$U$，使所考虑的\term{sheaf}{层}在$U$上局部常值，令$F$为其闭补集，记$U\xhookrightarrow{j}X\xhookleftarrow{i}F$。利用下述事实：给定$X$上的\term{sheaf}{层}，\term{equivalence}{等价}于给定$U$上的\term{sheaf}{层}$F_U$、$F$上的\term{sheaf}{层}$F_F$以及\term{recollement}{粘合}态射$F_F\to i^*j_*F_U$。经典拓扑和\term{etale}{平展}拓扑中都如此。归纳假设适用于$F$，所考虑的\term{sheaf}{层}在$U$上局部常值，又有$j_*$的比较定理，因此结论成立。

\SourcePage{158}
由(A)容易推出(A$'$)和(A$''$)。

通常的\term{cohomology}{上同调}运算$(f^*,Rf^!,Rf_*,Rf_!)$可以归约到$f$为嵌入时的$f_!$，以及任意$f$的$f^*$和$Rf_*$：

对$Rf_!$，写$f=gh$，其中$g$\term{proper}{固有}，$h$为开嵌入。于是$Rf_!=Rg_*h_!$。

对$Rf^!$，写$f=gh$，其中$g$\term{smooth}{光滑}，$h$为闭嵌入。于是$Rf^!=Rh^!Rg^!$。Poincaré\term{duality}{对偶}定理通过$g^*$计算$Rg^!$：当$g$的相对维数\term{pure}{纯}为$d$时，$Rg^!K=g^*K[2d](d)$。最后，设$j$为$h$的像的开补集的包含映射，利用三角$(h_*Rh^!,\allowbreak\id,\allowbreak Rj_*j^*)$便可用$Rj_*$控制$Rh^!$。这种$f=gh$的分解未必全局存在，但局部存在；一旦定义了比较箭头，这就足够了。比如，可通过$Rf_!$的比较同构以及\term{adjunction}{伴随}对$(Rf_!,Rf^!)$来定义比较箭头。

利用公式
\[
Rf_*\bigl(R\mathcal Hom(K,Rf^!L)\bigr)
 =R\mathcal Hom(f_!K,L),
\]
经过逐步分解，$R\mathcal Hom(K,L)$的情形可归约到$Rf_*$、$f$为嵌入时的$f_!$以及$K$为\term{local-system}{局部常值层}的情形。

\par\medskip\noindent\textbf{6.1.3.} \term{vanishing-cycles}{消失轮换}理论中也有比较定理，陈述见SGA~7 XIV 2.8。

\par\medskip\noindent\textbf{6.1.4.} 在忘掉$\C$的拓扑这第一步中，务必计入\term{tate-twist}{Tate扭转}：
\begin{itemize}
\item 对维数为$d$的$X$，基本类是态射$H_c^{2d}(X,\Z/\ell^n(d))\to\Z/\ell^n$。
\item 若$X$和$Y$均\term{smooth}{光滑}，且$X$在$Y$中的余维\term{pure}{纯}为$d$，则\term{gysin-morphism}{Gysin态射}为
\[
H^i(X,\Z/\ell^n)\longrightarrow H^{i+2d}(Y,\Z/\ell^n(d)),\quad\ldots
\]
\end{itemize}
以下几点有助于辨认这些扭转：
\begin{itemize}
\item 如果把$\C$看作$\R$的\emph{某个}代数闭包，而不指明哪一个平方根是$i$、哪一个是$-i$，那么这些扭转在经典理论中就已出现，至少模$2$如此。事实上，令$\Z(1)=\Ker(\exp:\C\to\C^*)=2\pi i\Z$。选择$i$便通过$1\mapsto2\pi i$把$\Z$与$\Z(1)$等同；将$i$换成$-i$，这个同构就变为其负。同样，选择$i$确定$\C$的一个定向，进而确定每个\term{smooth}{光滑}簇的定向；
\SourcePage{159}
把$i$换成$-i$，就会反转奇数维簇的定向。因此，不依赖于$i$的选择，\term{pure}{纯}维数为$d$的\term{smooth}{光滑}簇$X$的定向\term{sheaf}{层}可与取值为$\Z(d):=\Z(1)^{\otimes d}$的\term{constant-sheaf}{常值层}等同；基本类将$H_c^{2d}(X,\Z(d))$映到$\Z$。\term{l-adic}{$\ell$-进}基本类由这个整系数基本类，以及指数映射给出的同构得到：
\[
\Z(1)\otimes\Z/n\xrightarrow{\sim}\Z/n(1),\qquad z\longmapsto\exp(z/n).
\]
\item 在de Rham\term{cohomology}{上同调}中，它们对应于因子$(2\pi i)^d$。
\item 记忆方法：在\term{derived-category}{导出范畴}中，扭转$(d)$往往\term{adjunction}{伴随}着平移$[2d]$，例如$f$\term{smooth}{光滑}、相对维数为$d$时，$Rf^!K=f^*K[2d](d)$。注意，运算$K\mapsto K[2d](d)$保持\term{weight}{权}。
\end{itemize}

\par\medskip\noindent\textbf{6.1.5.} 前述工具并不总能直接套用。例如，可以用奇异链定义交叉（上）同调：要求链及其边界与适当\term{stratification}{分层}中各\term{stratum}{层片}的交的维数不能过大。这个定义本身不便转换成\term{etale}{平展}\term{cohomology}{上同调}的语言。但是，一旦将\term{intersection-cohomology}{交叉上同调}等同于\term{hypercohomology}{超上同调}$\mathbb H^*(X,j_{!*}(\Q[d]))$，其中$j$为某个维数为$d$的\term{smooth}{光滑}稠密开集的包含映射，而$j_{!*}$由2.1.11描述，转换就立即完成。

\par\medskip\noindent\textbf{6.1.6.} 接下来将讨论第二步：对不涉及$\C$的拓扑的命题，从$\C$过渡到$\overline{\F}$。在此之前，尽管逻辑上并非必需，我们先举例说明如何从一个$\C$上的簇过渡到一个数域上定义的簇。

设$X$为复射影簇。它在$\mathbb P^n(\C)$中由有限多个齐次方程$P_\alpha=0$定义。记$P_\alpha$的系数为$a_{\alpha\beta}$。如果把$a_{\alpha\beta}$视为不定元，那么这些$P_\alpha$就定义了一个以\term{affine}{仿射}空间$S$为参数空间的射影簇族，从而得到射影态射$f:Y\to S$；$X$正是$S$中坐标等于各$a_{\alpha\beta}$原来取值的点$s$处的纤维。这个态射$f:Y\to S$由某个$\Q$上\term{scheme}{概形}态射经标量从$\Q$\term{extension}{扩张}至$\C$得到。
\SourcePage{160}
我们就说，它定义在$\Q$上。

设$f:Y\to S$为$\Q$上\term{finite-type}{有限型}\term{scheme}{概形}之间的分离态射。存在$S$的一个定义在$\Q$上的\term{stratification}{分层}$\mathcal T$，使得对每个\term{stratum}{层片}$T$，由$f$限制到$T$所得的态射$f_T:Y_T\to T$的纤维在“拓扑意义下局部常值”。例如，$Y_T(\C)\to T(\C)$是一个拓扑上局部平凡的纤维化。这个结论来自Thom同痕定理，以及Whitney条件可用代数方式表述这一事实（见J.-L. Verdier，\emph{Inventiones} \textbf{36} (1976)，295--312页，推论5.1）。

每个\term{stratum}{层片}的每个连通分支都有定义在某个数域上的点。因此，若$s\in S(\C)$，而$t$与$s$位于同一\term{stratum}{层片}的同一连通分支，且$t$定义在某个数域上，则$Y_s$与$Y_t$同胚。比如，已知$Y_t$的\term{intersection-cohomology}{交叉上同调}满足\term{hard-lefschetz}{强Lefschetz定理}，要由此证明复射影簇$X=Y_s$的相同结论：带有$\mathcal O(1)$的第一\term{chern-class}{Chern类}$\ell$的$X$，与$(Y_t,\ell)$同胚，于是结论成立。这个论证用到了\term{intersection-cohomology}{交叉上同调}的拓扑不变性（[6]）。如果在带有适当\term{stratification}{分层}的$Y$上工作，也可以不用这一不变性。

我们往往关心的并非单个簇$X$，而是态射$u:X_1\to X_2$。此时同样可以找到定义在$\Q$上的图
\[
\begin{tikzcd}[ampersand replacement=\&]
Y_1\arrow[rr,"{u}"]\arrow[dr,"{f_1}"']\&\&Y_2\arrow[dl,"{f_2}"]\\
\&S\&
\end{tikzcd}
\]
使$u$由它在$s\in S(\C)$处取纤维得到。需要留意：一般不存在$S$的\term{stratification}{分层}$\mathcal T$，使$u_t:f_1^{-1}(t)\to f_2^{-1}(t)$的拓扑类型在同一\term{stratum}{层片}的每个连通分支上保持不变；其拓扑类型可能连续变化。不过，仍然存在定义在$\Q$上的$S$的\term{stratification}{分层}$\mathcal T$，使得从$S$到$T$作\term{base-change}{基变换}后，态射
\SourcePage{161}
\[
\begin{tikzcd}[ampersand replacement=\&]
Y_{1T}\arrow[rr,"{u_T}"]\arrow[dr]\&\&Y_{2T}\arrow[dl]\\
\&T\&
\end{tikzcd}
\]
具有下列性质：任意$T'\to T$给出笛卡尔图
\[
\begin{tikzcd}[ampersand replacement=\&]
Y_{1T'}(\C)\arrow[r,"{v}"]\arrow[d,"{u}"']\&Y_{1T}(\C)\arrow[d,"{u}"]\\
Y_{2T'}(\C)\arrow[r,"{v}"']\&Y_{2T}(\C),
\end{tikzcd}
\]
且对$Y_{1T}(\C)$上的任意\term{local-system}{局部系统}$L$，有
\[
v^*R^qu_*L\xrightarrow{\sim}R^qu_*v^*L.
\]
适当推广这个结果（参见6.1.2中的归约）后，实际应用中就已足够。

\par\medskip\noindent\textbf{6.1.7.} EGA~IV §§8、9列出了许多如下类型的结果。设交换环$A$为环$A_i$的归纳极限；为简便起见，假设$A$与各个$A_i$均为\term{noetherian}{诺特}环。$\Spec(A)$上任何“有限表示”的对象$X$（例如\term{finite-type}{有限型}\term{scheme}{概形}、这类\term{scheme}{概形}间的态射、这类\term{scheme}{概形}上的\term{coherent-sheaf}{凝聚层}等），都来自某个$\Spec(A_i)$上的同类对象$X_i$；对充分大的$j$，两种$X_i$的选择给出相同的$X_j:=X_i\times_{\Spec A_i}\Spec A_j$。设某个“相对”性质$P$在\term{base-change}{基变换}下稳定，则$X/\Spec A$具有$P$，当且仅当充分大的$j$对应的$X_j/\Spec A_j$具有$P$。

这里我们关心的情形是$A=\C$，将它写成自身各个有限生成$\Z$-子代数的归纳极限。

\noindent\textbf{变体。} 对$\C$中任意有限生成$\Z$-子代数$A$，存在$A$中某个$f\ne0$，使$\Spec(A[1/f])$在$\Spec\Z$上\term{smooth}{光滑}。用一个共尾子集$I$替换$\C$中所有有限生成$\Z$-子代数构成的\term{filtration}{滤过}有序集，仍有$\C=\varinjlim A_i$，而且可使这个归纳系中的每个$S_i:=\Spec(A_i)$都在$\Spec\Z$上\term{smooth}{光滑}。

\SourcePage{162}
于是，$\C$上任意\term{finite-type}{有限型}\term{scheme}{概形}$X$都来自某个$X_i/S_i$；若$X,Y$来自$X_i,Y_i/S_i$，则态射（分别地，同构）$f:X\to Y$来自某个$f_j:X_j\to Y_j$，其中$j\ge i$，这里记$X_j=X_i\times_{S_i}S_j$，$Y$同理。而$f_j$的两种选择，在某个$X_k$上就相等（$k\ge j$）。因此，可以说“当$i$充分大时，$X_i$本质唯一”。

同样，$X$的\term{stratification}{分层}$\mathcal T$、\term{constructible}{可构造}$\Z/\ell^n$-模\term{sheaf}{层}、这种\term{sheaf}{层}之间的态射等，都来自某个充分大的$i$对应的$X_i$上的同类对象；$X_i$上的两种选择，在适当的$X_j$上便相等，其中$j\ge i$。

\noindent\textbf{记号。} 对$\C$上的一个“对象”$\Xi$，用$\Xi_S$表示$S=\Spec(A)$上的一个“对象”，其中$A\subset\C$是在$\Z$上有限生成且取得足够大的子环，使$\Xi$由$\Xi_S$经$A$到$\C$的标量\term{extension}{扩张}得到。对$S$上的$T$，用$\Xi_T$表示$\Xi_S$在$T$上的拉回。这个记号有歧义（一般说来，$\Xi_S$并不唯一），但上述结果保证这种歧义不会造成妨碍。

若$\mathcal T$为$X$的一个\term{stratum}{层片}\term{smooth}{光滑}（分别地，连通）的\term{stratification}{分层}，那么把$S$取得更大之后，$\mathcal T_S$的各\term{stratum}{层片}就在$S$上\term{smooth}{光滑}（分别地，具有连通几何纤维）。若$X$上的\term{sheaf}{层}$F$在各\term{stratum}{层片}上局部常值，那么把$S$取得更大后，$X_S$上的$F_S$也具有这一性质。

设$f:X\to Y$为$\C$上\term{finite-type}{有限型}\term{scheme}{概形}之间的态射，$F$为$X$上的\term{constructible}{可构造}$\Z/\ell^n$-模\term{sheaf}{层}。如上取$S=\Spec(A)$上的$f_S:X_S\to Y_S$及$X_S$上的$F_S$。由SGA~$4\frac12$（有限性定理1.9），存在非空开集$U\subset S$，使任意\term{scheme}{概形}态射$g:U'\to U$给出的笛卡尔图
\[
\begin{tikzcd}[ampersand replacement=\&]
X_{U'}\arrow[r,"{g}"]\arrow[d,"{f}"']\&X_U\arrow[d,"{f}"]\\
Y_{U'}\arrow[r,"{g}"']\&Y_U
\end{tikzcd}
\]
都满足
\[
g^*R^qf_*F_U\xrightarrow{\sim}R^qf_*g^*F_U,
\]
且各个$R^qf_*F_U$均\term{constructible}{可构造}。

\noindent\textbf{推论}（参见6.1.2中的归约）。给定有限多个$\C$上的\term{finite-type}{有限型}\term{scheme}{概形}$X_i$、它们之间有限多个态射，以及这些$X_i$上有限多个\term{constructible}{可构造}$\Z/\ell^n$-模\term{sheaf}{层}，
\SourcePage{163}
可以把它们全部下降到$S=\Spec(A)$，其中$A\subset\C$是在$\Z$上有限生成且取得足够大的子环。适当缩小$S$后，可以假设，对给定的\term{constructible}{可构造}\term{sheaf}{层}作有界次数的通常\term{cohomology}{上同调}运算
\[
R^*f_*,\ R^*f_!,\ f^*,\ R^*f^!,\ \otimes,\ \mathcal Tor_p,\ \mathcal Ext^p
\]
所得的\term{sheaf}{层}仍然\term{constructible}{可构造}，并与任意\term{base-change}{基变换}$S'\to S$交换。如果$a_i$为$X_{iS}$到$S$的投影，还可假设，对这些运算的结果应用$R^*a_{i*}$所得的\term{sheaf}{层}在$S$上局部常值。

\par\medskip\noindent\textbf{6.1.8.} 上面讨论的是$\Z/\ell^n$-模\term{sheaf}{层}。$\Z_\ell$-\term{sheaf}{层}本质上是无限性对象，一般不能从$\C$下降到$S$。例如，$\mathbb G_{m\C}=\mathbb P^1(\C)-\{0,\infty\}$上存在\term{monodromy}{单值作用}并非拟幺幂的\term{smooth}{光滑}$\Z_\ell$-\term{sheaf}{层}。这样的$\Z_\ell$-\term{sheaf}{层}不来自$\mathbb G_{mK}$上的$\Z_\ell$-\term{sheaf}{层}，其中$K\subset\C$为$\Q$的有限生成\term{extension}{扩张}。

下面给出绕过这一障碍的方法。设$X$在$\C$上\term{finite-type}{有限型}，$\mathcal T$是$X$的一个\term{stratification}{分层}（2.1），各\term{stratum}{层片}\term{smooth}{光滑}且连通；数据$\mathcal L$为每个\term{stratum}{层片}$T\in\mathcal T$指定有限族$\mathcal L(T)$，其成员为不可约的\term{smooth}{光滑}$\Z/\ell$-模\term{sheaf}{层}。设$(X_S,\mathcal T_S,\mathcal L_S)$定义在$S=\Spec(A)$上，$A\subset\C$在$\Z$上有限生成，且经$A$到$\C$的标量\term{extension}{扩张}给出$(X,\mathcal T,\mathcal L)$。假设$S$已取得足够大，使$\mathcal T_S$中各\term{stratum}{层片}$T$都在$S$上\term{smooth}{光滑}且几何纤维连通；还要求，对于形如$j_!L$的$F,G$，其中$j:T\hookrightarrow X_S$为$\mathcal T_S$中的\term{stratum}{层片}，$L\in\mathcal L_S(T)$，各个$\mathcal Ext^q(F,G)$都与\term{base-change}{基变换}$S'\to S$相容，而且若$a$为到$S$的投影，则各个$R^pa_*\mathcal Ext^q(F,G)$局部常值并与\term{base-change}{基变换}相容。

取$A\subset V\subset\C$，其中$V$为严格Hensel离散赋值环，其剩余域是$A$的某个有限剩余域的代数闭包。对$\Spec A$的任意闭点$a$，都存在这样的$V$，使态射$\Spec V\to\Spec A$将$\Spec V$的闭点映到$a$：在完备化$\widehat A_a$中取一条“解析弧”，即一个一维商，使其不包含在$\Spec A$任何非平凡闭子集在$\Spec\widehat A_a$中的\term{inverse-image}{逆像}内，再作正规化和严格Hensel化即可。令$(X_V,\mathcal T_V,\mathcal L_V)$由$(X_S,\mathcal T_S,\mathcal L_S)$\term{base-change}{基变换}得到，$(X_s,\mathcal T_s,\mathcal L_s)$为其特殊纤维。考虑\term{scheme}{概形}图
\SourcePage{164}
\[
X\xrightarrow{u}X_V\xleftarrow{i}X_s.
\]
记$D_{\mathcal T,\mathcal L}^b(X,\Z_\ell)$为$D_c^b(X,\Z_\ell)$中由满足下述条件的$K$组成的子范畴：各个$H^iK$模$\ell$约化后，限制到任意\term{stratum}{层片}$T$，都是\term{local-system}{局部常值层}，且其组成因子属于$\mathcal L(T)$。类似地定义$D_{\mathcal T,\mathcal L}^b(X_V,\Z_\ell)$和$D_{\mathcal T,\mathcal L}^b(X_s,\Z_\ell)$。

\par\medskip\noindent\textbf{引理6.1.9.} \emph{沿用上述记号，函子}
\[
D_{\mathcal T,\mathcal L}^b(X,\Z_\ell)
\xleftarrow{u^*}D_{\mathcal T,\mathcal L}^b(X_V,\Z_\ell)
\xleftarrow{i^*}D_{\mathcal T,\mathcal L}^b(X_s,\Z_\ell)
\]
\emph{是范畴\term{equivalence}{等价}。}\footnote{译注：原书将第二个标为$i^*$的箭头印为向左；依照\term{inverse-image}{逆像}函子的定义，其方向应为从$X_V$对应的范畴指向$X_s$对应的范畴。这里保留原图方向以便对照。}

将$\Z_\ell$换成$\Z/\ell^n$的类似命题，形式地来自以下事实：对6.1.8中的$F,G$，各个$H^p(-,\mathcal Ext^q(F,G))$在$X,X_V,X_s$上“相同”。可以检查有限Tor维数得到保持。因此，按照\term{l-adic}{$\ell$-进}\term{derived-category}{导出范畴}的定义，将$D_{\mathcal T,\mathcal L}^b(-,\Z_\ell)$换成$D_{\mathcal T,\mathcal L,t.f.}^b(-,\Z/\ell^n)$（t.f.表示有限Tor维数）所得到的命题，就推出6.1.9。

\par\medskip\noindent\textbf{6.1.10.} 因而，对给定的$(X,\mathcal T,\mathcal L)$，以及任意足够大的$S=\Spec(A)$和任意上述$V$（$A\subset V\subset\C$），都得到一个范畴\term{equivalence}{等价}
\begin{equation*}
D_{\mathcal T,\mathcal L}^b(X,\Z_\ell)
\longleftrightarrow D_{\mathcal T,\mathcal L}^b(X_s,\Z_\ell),
\tag{6.1.10.1}
\end{equation*}
它把$\C$与一个有限域的代数闭包联系起来。

增大$S$之后，两种$\mathcal T_S,\mathcal L_S$的选择会相等。这一事实也可如下表述。设$\mathcal V$为$\C$中所有剩余域为某个有限域的代数闭包的离散赋值环$V$组成的集合。对在$\Z$上有限生成的$A\subset\C$，记$\mathcal V(A)$为满足$A\subset V$的$V$的集合。这些$\mathcal V(A)$构成$\mathcal V$的子集的一个\term{filtration}{滤过}有序集，即一个滤子基；我们将考虑沿相应滤子的芽。

设$X$为$\C$上\term{finite-type}{有限型}\term{scheme}{概形}，由$X_S$经标量\term{extension}{扩张}得到，其中$S=\Spec(A)$。则由$X_s$组成的系的芽，仅依赖于$X$，且在唯一同构意义下如此；这里$s$为$\Spec(V)$的闭点，$V\in\mathcal V(A)$。其他所考虑的“有限表示”对象也一样。
\SourcePage{165}
同样，对给定的$(X,\mathcal T,\mathcal L)$，范畴\term{equivalence}{等价}系(6.1.10.1)的芽仅依赖于$(X,\mathcal T,\mathcal L)$。

若$\mathcal T',\mathcal L'$加细$\mathcal T,\mathcal L$，则$\mathcal T,\mathcal L$所对应的\term{equivalence}{等价}系(6.1.10.1)的芽，当然由$(\mathcal T',\mathcal L')$的类似芽诱导。但是，因为我们只有芽，不能取极限来得到\term{equivalence}{等价}$D_c^b(X,\Z_\ell)\longleftrightarrow D_c^b(X_s,\Z_\ell)$。只能断言：每个$K\in D_c^b(X,\Z_\ell)$确定一个对象系的芽，其中各$K_s$属于$D_c^b(X_s,\Z_\ell)$。

若$(\mathcal T,\mathcal L)$满足：(a)各\term{stratum}{层片}\term{smooth}{光滑}且连通；(b)对\term{stratum}{层片}$j:T\hookrightarrow X$、$L\in\mathcal L(T)$及$T'\in\mathcal T$，各个$R^qj_*L|T'$均局部常值，组成因子属于$\mathcal L(T')$；则可取$S$使同样的性质在$S$上成立，并使各个$R^qj_*L$与任意\term{base-change}{基变换}$S'\to S$交换。此时，对$\mathcal T$上的任意\term{perversity}{反常度}函数$p$，6.1.9中的\term{equivalence}{等价}都保持\term{perversity}{反常度}为$p$的\term{t-structure}{$t$-结构}，\term{equivalence}{等价}系(6.1.10.1)的芽也一样。

利用6.1.7，可以使6.1.9中的\term{equivalence}{等价}（以及\term{equivalence}{等价}系(6.1.10.1)的芽）保持通常的\term{cohomology}{上同调}运算。例如，对$\mathrm T=R\mathcal Hom$或$\overset L\otimes$，可以把每个$(\mathcal T,\mathcal L)$扩大到某个$(\mathcal T',\mathcal L')$，使$\mathrm T$将$D_{\mathcal T,\mathcal L}^b\times D_{\mathcal T,\mathcal L}^b$映入$D_{\mathcal T',\mathcal L'}^b$，然后取$S$，使上述$V$对应的6.1.9中的\term{equivalence}{等价}保持函子$\mathrm T$。同样，对$f:X\to Y$（或$Y\to X$），可在$Y$上找到$(\mathcal T',\mathcal L')$，使$\mathrm T=Rf_*,Rf_!$（或$f^*,Rf^!$）将$D_{\mathcal T,\mathcal L}^b$映入$D_{\mathcal T',\mathcal L'}^b$，并且6.1.9中的\term{equivalence}{等价}与函子$\mathrm T$交换。

对$R\mathcal Hom$和$\overset L\otimes$，还可利用$\mathcal L(T)$是$\Z/\ell$-模\term{sheaf}{层}的同构类集合，且这些\term{sheaf}{层}的\term{monodromy}{单值作用}有限这一事实，用更简单的方法把$(\mathcal T,\mathcal L)$加细为$(\mathcal T',\mathcal L')$，使$D_{\mathcal T',\mathcal L'}^b$在$R\mathcal Hom$和$\overset L\otimes$下稳定。

缩小$S$之后，还可以假设(6.1.10.1)限制到\term{sheaf}{层}上时保持下列性质：
\begin{itemize}
\item $F$在同一\term{stratum}{层片}各点处的纤维的$\Z_\ell$-模结构，即秩及\term{torsion}{挠}部分的初等因子。
\SourcePage{166}
\item $F$的\term{smooth}{光滑}性。注意到$\Z_\ell$-\term{sheaf}{层}\term{smooth}{光滑}，当且仅当$\ell:F\to F$的核与余核均\term{smooth}{光滑}，就可归约到$\Z/\ell$-模\term{sheaf}{层}的情形。
\item 当$X$\term{smooth}{光滑}连通、$F$在$X$上\term{smooth}{光滑}时，保持其\term{monodromy}{单值作用}。先缩小$X$，使它只有一个\term{stratum}{层片}，这不改变单值群；再取一个连通有限覆盖，使$\mathcal L(X)$中的元素都平凡化。然后利用：$X\to X_V$和$X_s\to X_V$在$H^1(-,\Z/\ell)$上诱导同构，因此它们在$\pi_1$的pro-$\ell$完备化上诱导满射。
\end{itemize}
此外，由$\Z_\ell$-\term{derived-category}{导出范畴}的\term{equivalence}{等价}系(6.1.10.1)，立即可得$\Q_\ell$、$E_\lambda$和$\overline{\Q}_\ell$-\term{derived-category}{导出范畴}的类似\term{equivalence}{等价}系。要在$\Q_\ell$-\term{sheaf}{层}的情形验证它保持\term{smooth}{光滑}性，可利用下列事实：
\begin{enumerate}[label=(\alph*)]
\item $F$\term{smooth}{光滑}，当且仅当它在$X_{\mathrm{red}}$的正规化上的\term{inverse-image}{逆像}\term{smooth}{光滑}。
\item 设$X$正规且连通，$j:U\hookrightarrow X$为稠密开集。$X$上的$F$\term{smooth}{光滑}，当且仅当其所有纤维维数相同，它在$U$上的限制\term{smooth}{光滑}，并且$F\xrightarrow{\sim}j_*j^*F$。
\end{enumerate}

\section*{6.2. 例子}
\addcontentsline{toc}{section}{6.2. 例子}
\par\medskip\noindent\textbf{6.2.1. \term{affine}{仿射}\term{scheme}{概形}的\term{cohomology}{上同调}维数。}

设$X$为$\C$上维数为$d$的\term{affine}{仿射}\term{scheme}{概形}。我们要从\term{l-adic}{$\ell$-进}\term{cohomology}{上同调}中的类似结果（4.1.4），推出：对$X(\C)$上任意\term{constructible}{可构造}$\C$-向量空间\term{sheaf}{层}$F$，当$i>d$时，$H^i(X(\C),F)=0$。以下简记$X(\C)$为$X$。逐步分解后，可假设$F$形如$j_!L$，其中$j:Y\hookrightarrow X$为连通子\term{scheme}{概形}的包含映射，$L$为$Y$上的复\term{local-system}{局部系统}。取$y\in Y$。\term{local-system}{局部系统}$L$由$\pi_1(Y,y)$在$L_y\simeq\C^n$上的一个表示定义。因为$\pi_1$有限生成，这个表示经由
\[
\rho_A:\pi_1(Y,y)\longrightarrow\mathrm{GL}(n,A)
\]
分解，其中$A\subset\C$在$\Z$上有限生成。于是$L=L_A\otimes_A\C$，其中$L_A$为某个$A$-模\term{local-system}{局部系统}。令$F_A=j_!L_A$。这里承认如下已知事实：各个$H^i(X,F_A)$为有限生成$A$-模，且对任意$A$-模$M$，有
\SourcePage{167}
\[
R\Gamma(X,F_A\otimes_A M)\xleftarrow{\sim}R\Gamma(X,F_A)\overset L\otimes_A M.
\]
这些有限性结果来自如下事实：若$\overline X$为$X$的射影完备化，则$(\overline X,\overline X-X)$可三角剖分。设$N$为使$H^i(X,F_A)\ne0$的最大整数$i$，则
\[
H^N(X,F_A\otimes_A M)=H^N(X,F_A)\otimes_A M.
\]
用反证法证明$N\le d$。假设$N>d$，并取$M=A/\mathfrak m$，其中$\mathfrak m$为$A$的极大理想。记$F_{A/\mathfrak m}=F_A\otimes_A A/\mathfrak m$。由于$A/\mathfrak m$有限，$F_{A/\mathfrak m}$为$X_{\mathrm{\acute et}}$上某个\term{sheaf}{层}$F_{A/\mathfrak m}^{\mathrm{\acute et}}$的$\varepsilon^*$。于是
\[
H^N(X,F_A)\otimes_A A/\mathfrak m
=H^N(X,F_{A/\mathfrak m})
\xleftarrow[\text{6.1.2(C)}]{\sim}
H^N(X_{\mathrm{\acute et}},F_{A/\mathfrak m}^{\mathrm{\acute et}})=0,
\]
最后一个等号是因为$N>d$。这对每个$\mathfrak m$都成立，而$H^N(X,F_A)$有限生成，故$H^N(X,F_A)=0$，矛盾。因此，各个$H^i(X,F_A)$在$i>d$时均为零；再利用$H^i(X,F)=H^i(X,F_A)\otimes_A\C$，就得到结论。

\noindent\textbf{变体。} 不先逐步分解$F$以归约到$F=j_!L$，也可以直接写成$F=F_A\otimes_A\C$，其中$A$如上，$F_A$为纤维是自由$A$-模的\term{constructible}{可构造}\term{sheaf}{层}。

\noindent\textbf{注。} 所得结果还不够理想：只要$X$为Stein空间，$F$为某个局部有限解析\term{stratification}{分层}的各\term{stratum}{层片}上局部常值的阿贝尔群\term{sheaf}{层}，就有$H^i(X,F)=0$（$i>\dim_\C X$）。这可用Morse理论的适当变体证明。

\par\medskip\noindent\textbf{6.2.2. \term{weight-filtration}{权滤过}。}

设$X$定义在$\C$上。为了得到$X$的\term{cohomology}{上同调}上的\term{mixed}{混合}Hodge结构，首先用一个带有适当紧化的单\term{pure}{纯}\term{scheme}{概形}$X_*\hookrightarrow\overline X_*$替换$X$（[3]）。单\term{pure}{纯}\term{scheme}{概形}本质上是无限性对象：它包含无穷多个$X_i$和无穷多个箭头。因此，未必能在$S=\Spec(A)$上找到一个版本，其中$A\subset\C$在$\Z$上有限生成。不过，要研究$i\le n$时的$H^i(X)$，只需要$i\le n+1$时的$X_i$和$\overline X_i$；这个\term{truncation}{截断}
\[
(X_i)_{i\le n+1}\hookrightarrow(\overline X_i)_{i\le n+1}
\]
可下降到$S$。$H^*(X(\C),\Q)$的\term{weight-filtration}{权滤过}$W_\Q$，是由$X_*\hookrightarrow\overline X_*$构造的某个\term{spectral-sequence}{谱序列}的终端\term{filtration}{滤过}。同一\term{spectral-sequence}{谱序列}也存在于\term{l-adic}{$\ell$-进}\term{cohomology}{上同调}中。因此，$H^n(X,\Q_\ell)=H^n(X(\C),\Q)\otimes\Q_\ell$上由$W_\Q$得到的\term{filtration}{滤过}$W_\ell$也能保留到$S$上：
\SourcePage{168}
对$a:X_S\to S$，存在$R^na_*\Q_\ell$上的一个\term{filtration}{滤过}$W_\ell$，它在$\C$上给出$W_\ell$。缩小$S$之后，可以假设，$\C$上对$X_i$和$\overline X_i$（$i\le n+1$）施加的条件在$S$上仍成立：$\overline X_{iS}$在$S$上\term{proper}{固有}且\term{smooth}{光滑}，$X_{iS}\subset\overline X_{iS}$是某个相对正规交叉除子的补集，且
\[
X_{iS}\longrightarrow
\bigl(\operatorname{cosq}_{i-1}\operatorname{sq}_{i-1}(X_{*S})\bigr)_i
\]
满射且\term{proper}{固有}；同时，$W_\ell$作为该\term{spectral-sequence}{谱序列}终端\term{filtration}{滤过}的构造与任意\term{base-change}{基变换}交换。利用Weil猜想可见，在$S$的每个闭点$s$处，$W_\ell$诱导$H^n(X_{\bar s},\Q_\ell)$的\term{weight-filtration}{权滤过}，即由\term{frobenius}{Frobenius}特征值的复绝对值所作的算术定义给出的\term{filtration}{滤过}。下面给出一个应用。

\par\medskip\noindent\textbf{命题6.2.3.} \emph{设$f:X\to Y$为$\C$上\term{finite-type}{有限型}\term{scheme}{概形}之间的分离态射。假设$R^nf_*\Q$局部常值，而且对$Y$中每个$y$，都有$(R^nf_*\Q)_y\xrightarrow{\sim}H^n(X_y,\Q)$。那么，各个$H^n(X_y,\Q)$的\term{weight-filtration}{权滤过}构成\term{local-system}{局部系统}$R^nf_*\Q$的一个局部常值\term{filtration}{滤过}。}

只需在张量$\Q_\ell$之后证明6.2.3。此时目标是从定理[1]~3.4.1推出结论：该定理断言，若$F_0$为$\F_q$上\term{finite-type}{有限型}\term{scheme}{概形}$Y_0$上的\term{smooth}{光滑}局部常值$\Q_\ell$-\term{sheaf}{层}，且$F_0$\term{mixed}{混合}，则$F_0$有一个由\term{smooth}{光滑}子\term{sheaf}{层}构成的\term{filtration}{滤过}$W$，使$\gr_i^WF_0$\term{pointwise-pure}{逐点纯}，\term{weight}{权}为$i$。

6.2.3中的困难在于：为了计算$X_y=f^{-1}(y)$的\term{cohomology}{上同调}的\term{weight-filtration}{权滤过}而构造的紧化单\term{pure}{纯}\term{scheme}{概形}$(X_y)_*$，并不总能对$y$一致地构造。不过，如果只考虑\term{truncation}{截断}单\term{pure}{纯}\term{scheme}{概形}，就存在$Y$的一个\term{stratification}{分层}，使在每个\term{stratum}{层片}$T$上方都可一致地构造；也就是说，存在增广到$X_T$的紧化\term{truncation}{截断}单\term{pure}{纯}\term{scheme}{概形}$X_{*T}\hookrightarrow\overline X_{*T}$，满足施加于$(X_y)_*$的条件的相对版本。

将$f:X\to Y$下降为$S=\Spec(A)$上的$f_S:X_S\to Y_S$，其中$A\subset\C$在$\Z$上有限生成。缩小$S$之后，“$R^nf_*\Q_\ell$与取纤维交换”这一假设由$f$传到$f_S$。事实上，存在$Y_S$的\term{stratification}{分层}$\mathcal T_S$，使对其中每个$Y'_S$，态射$f'_S:X'_S=f_S^{-1}(Y'_S)\to Y'_S$满足“$Rf'_{S*}\Q_\ell$的形成与任意\term{base-change}{基变换}交换”。
\SourcePage{169}
于是$f_S$所需的性质\term{equivalence}{等价}于：对$\mathcal T_S$中的每个$Y'_S$，笛卡尔图
\[
\begin{tikzcd}[ampersand replacement=\&]
X'_S\arrow[r,hook]\arrow[d,"{f'_S}"']\&X_S\arrow[d,"{f_S}"]\\
Y'_S\arrow[r,hook,"{u}"']\&Y_S
\end{tikzcd}
\]
给出同构$u^*R^nf_{S*}\Q_\ell\xrightarrow{\sim}R^nf'_{S*}\Q_\ell$。这个态射在$\C$上是同构，所以在$S$的某个非空开集上方也是同构；把$S$缩小到这个开集即可。

综合这些论证，得到$Y$的一个\term{stratification}{分层}$\mathcal T$，在下降到适当的$S$后，对每个$Y'_S\in\mathcal T_S$，态射$f'_S:X'_S\to Y'_S$都满足：
\begin{enumerate}[label=\alph*)]
\item $R^nf'_{S*}\Q_\ell$与取纤维交换；
\item $R^nf'_{S*}\Q_\ell$为\term{smooth}{光滑}\term{sheaf}{层}，且有一个由\term{smooth}{光滑}子\term{sheaf}{层}构成的\term{weight-filtration}{权滤过}$W$，它在每个$H^n(X_y,\Q)\otimes\Q_\ell$上（$y\in Y'(\C)$）诱导其\term{weight-filtration}{权滤过}，并且在$S$的每个闭点上方诱导\term{l-adic}{$\ell$-进}\term{weight-filtration}{权滤过}。
\end{enumerate}
最后只需注意：对适当的$S$，6.1.10中的\term{equivalence}{等价}具有如下性质。若$Y$上的\term{smooth}{光滑}\term{sheaf}{层}$F$在某个\term{stratification}{分层}$\mathcal T$的每个\term{stratum}{层片}$T$上都带有\term{filtration}{滤过}$W_T$，且$F,W_T$分别对应于$F',W'_T$，则各个$W_T$能\term{recollement}{粘合}成$F$的一个由\term{smooth}{光滑}子\term{sheaf}{层}构成的\term{filtration}{滤过}，当且仅当$F'$和各个$W'_T$也能如此。这是因为，各个$W_T$能够\term{recollement}{粘合}\term{equivalence}{等价}于\term{recollement}{粘合}箭头保持$W$；\term{smooth}{光滑}性则参见6.1.10。

\noindent\textbf{注。} 除了这个算术证明，我们不知道6.2.3的其他证明。

\par\medskip\noindent\textbf{6.2.4.} 设$F$为$X(\C)$上单的$\C$-向量空间\term{perverse-sheaf}{反常层}，其中$X$为$\C$上\term{finite-type}{有限型}\term{scheme}{概形}。如果$F$属于满足下列条件的最小集合，就称$F$具有\term{geometric-origin}{几何来源}：
\begin{enumerate}[label=(\alph*)]
\item 它包含$X$为一点时的\term{constant-sheaf}{常值层}$\C$，并且在以下运算下稳定；
\SourcePage{170}
\item 对\term{scheme}{概形}态射$f$，取${}^pH^i\mathrm T$的组成因子，其中$\mathrm T=Rf_*,Rf_!,Rf^*,Rf^!$；
\item 取${}^pH^i(A\overset L\otimes B)$或${}^pH^iR\mathcal Hom(A,B)$的组成因子。
\end{enumerate}
在各个$D_c^b(X,\C)$中，取满足下述条件的$K$组成的子范畴：各个${}^pH^i(K)$的单\term{subquotient}{子商}都具有\term{geometric-origin}{几何来源}。这些子范畴在$f_*,f_!,f^*,f^!,\otimes$和$R\mathcal Hom$下稳定。由此还可得到其他稳定性。\footnote{译注：原书本句末的“参见”之后未印出引文。}还可证明它们在“\term{vanishing-cycles}{消失轮换}”构造下稳定。

若$X(\C)$上的\term{perverse-sheaf}{反常层}$F$是具有\term{geometric-origin}{几何来源}的单\term{perverse-sheaf}{反常层}的直和，就称它\term{semisimple}{半单}且具有\term{geometric-origin}{几何来源}。若$K\in D_c^b(X(\C),\C)$同构于各个$({}^pH^i(K))[-i]$的直和，且各个${}^pH^iK$均\term{semisimple}{半单}且具有\term{geometric-origin}{几何来源}，就称$K$\term{semisimple}{半单}且具有\term{geometric-origin}{几何来源}。

例如，设$j:U\hookrightarrow X$为\term{smooth}{光滑}连通、维数为$d$的Zariski开集，且$U$上的\term{local-system}{局部系统}$L$的\term{monodromy}{单值作用}有限。则$j_{!*}L[d]$\term{semisimple}{半单}且具有\term{geometric-origin}{几何来源}：因为$L$是某个有限覆盖$\pi:\widetilde X\to X$的$\pi_*\C$的直和因子，而$\widetilde X$上的\term{constant-sheaf}{常值层}$\C$是（一点）上的$\C$的\term{inverse-image}{逆像}。

\par\medskip\noindent\textbf{定理6.2.5（\term{decomposition-theorem}{分解定理}）。} \emph{设$f:X\to Y$为\term{proper}{固有}态射。如果$X(\C)$上的$K$\term{semisimple}{半单}且具有\term{geometric-origin}{几何来源}，则$Y(\C)$上的$f_*K$也如此。}

定理6.2.5\term{equivalence}{等价}于将系数域$\C$换成与之同构的域$\overline{\Q}_\ell$所得的版本，进而\term{equivalence}{等价}于再将经典拓扑换成\term{etale}{平展}拓扑所得的版本。

\par\medskip\noindent\textbf{引理6.2.6.} \emph{设$F$为单的、具有\term{geometric-origin}{几何来源}的$\overline{\Q}_\ell$-\term{perverse-sheaf}{反常层}。当$S=\Spec(A)$足够大时，6.1.10中的每个\term{equivalence}{等价}}
\[
D_{\mathcal T,\mathcal L}^b(X,\overline{\Q}_\ell)
\longleftrightarrow
D_{\mathcal T,\mathcal L}^b(X_s,\overline{\Q}_\ell)
\]
\emph{都把$F$变成$X_s$上一个单$\overline{\Q}_\ell$-\term{perverse-sheaf}{反常层}$F_s$，且满足下列条件：}

\noindent\textbf{(P)} \emph{$X_s$由某个$X_0/\F_q$经标量\term{extension}{扩张}得到，$F_s$由某个$F_0$经标量\term{extension}{扩张}得到，且$F_0$是\term{pure}{纯}的。}

\SourcePage{171}
$F_s$是单\term{perverse-sheaf}{反常层}这一点来自6.1.10。“具有\term{geometric-origin}{几何来源}”的定义是归纳性的，因此只需证明条件(P)在该定义涉及的各个运算下稳定。这些运算与所考虑的\term{equivalence}{等价}交换，剩下只需知道：
\begin{enumerate}[label=(\alph*)]
\item （一点）上的$\overline{\Q}_\ell$满足(P)。
\item 6.2.4中的运算${}^pH^i\mathrm T$将满足(P)的$F_s$变成由某个\term{mixed}{混合}$G_0$得到的$G_s$。$G_0$的单组成因子都是\term{pure}{纯}的（5.3.4），而由6.1.9和6.1.10，$G_s$的单组成因子满足(P)。
\item $\otimes$和$R\mathcal Hom$也一样。
\end{enumerate}

\par\medskip\noindent\textbf{6.2.7. 6.2.5的证明。} 只需处理$K$为\term{perverse-sheaf}{反常层}$F$的情形。应用6.2.6，得到$F_s$满足(P)。再应用5.4.5和5.3.8，得到$Rf_{s*}F_s$是它的各个${}^pH^i[-i]$的直和，且这些部分均\term{semisimple}{半单}。由6.1.10，$Rf_*F$也具有相同性质，6.2.5得证。

\par\medskip\noindent\textbf{推论6.2.8（整体\term{invariant-cycles}{不变轮换}定理）。} \emph{设$f:X\to Y$为\term{proper}{固有}态射（$X,Y$为$\C$上\term{finite-type}{有限型}\term{scheme}{概形}），$K\in D_c^b(X,\C)$。设$V$为$Y$的一个Zariski开集，使$H^if_*K$在$V$上局部常值。如果$K$\term{semisimple}{半单}且具有\term{geometric-origin}{几何来源}，则有满射}
\begin{equation*}
H^i(X,K)\twoheadrightarrow H^0(V,H^if_*K).
\tag{6.2.8.1}
\end{equation*}
\emph{特别地，若$V$连通，$y\in Y$，记$X_y=f^{-1}(y)$，则$\pi_1(V,y)$作用于$(H^if_*K)_y\xrightarrow{\sim}H^i(X_y,K)$，而(6.2.8.1)改写为}\footnote{译注：原文写$y\in Y$；此处要以$y$为$V$的基点，实际需取$y\in V$。}
\begin{equation*}
H^i(X,K)\twoheadrightarrow H^i(X_y,K)^{\pi_1(V,y)}.
\tag{6.2.8.2}
\end{equation*}

\par\medskip\noindent\textbf{推论6.2.9（局部\term{invariant-cycles}{不变轮换}定理）。} \emph{设$f:X\to Y$为\term{proper}{固有}态射，$K\in D_c^b(X,\C)$，$y\in Y$。设$V$为$Y$的一个Zariski开集，使$H^if_*K$在$V$上局部常值。对以$y$为中心的充分小的球$B$，记$j$为$B\cap V$到$B$的包含映射，则有}
\[
H^0(B\cap V,H^if_*K)=(j_*j^*H^if_*K)_y,
\]
\emph{以及}
\[
H^0(B,H^if_*K)\xrightarrow{\sim}(H^if_*K)_y=H^i(X_y,K).
\]
\SourcePage{172}
\emph{如果$K$\term{semisimple}{半单}且具有\term{geometric-origin}{几何来源}，则对这样的$B$，有}
\begin{equation*}
H^i(X_y,K)\twoheadrightarrow H^0(B\cap V,H^if_*K).
\tag{6.2.9.1}
\end{equation*}
\emph{若$B\cap V$连通，$z\in B\cap V$，还可改写为}
\begin{equation*}
H^i(X_y,K)\twoheadrightarrow H^i(X_z,K)^{\pi_1(B\cap V,z)}.
\tag{6.2.9.2}
\end{equation*}

6.2.9中所述的球$B$的存在性是已知的。这里用球只是为了叙述方便；也可以使用构成邻域基的其他邻域，例如用三角剖分定义的邻域。说明这一点之后，再考虑6.2.5，证明就与5.4.7--5.4.9相同。

\par\medskip\noindent\textbf{定理6.2.10（相对\term{hard-lefschetz}{强Lefschetz定理}）。} \emph{设$f:X\to Y$为射影态射，$\ell$为某个相对丰沛可逆\term{sheaf}{层}的第一\term{chern-class}{Chern类}，$F$为$X$上\term{semisimple}{半单}且具有\term{geometric-origin}{几何来源}的\term{perverse-sheaf}{反常层}。对$i\ge0$，有}
\[
\ell^i:{}^pH^{-i}f_*F\xrightarrow{\sim}{}^pH^if_*F(i).
\]
这个命题由6.2.6和5.4.10推出。一个有用的特例是$Y$只有一个点的情形；它给出\term{intersection-cohomology}{交叉上同调}的\term{hard-lefschetz}{强Lefschetz定理}。

\par\medskip\noindent\textbf{注6.2.11.} 我们希望在前述命题中，将“具有\term{geometric-origin}{几何来源}”替换为“形如$j_{!*}L[d]$，其中$j$为维数为$d$的\term{smooth}{光滑}连通$Y$的包含映射，$L$是某个可极化Hodge结构变分的底\term{sheaf}{层}\term{local-system}{局部系统}”。还希望将代数簇替换为紧Kähler簇。

若$f$为紧Kähler簇之间的态射，我们既不知道如何证明$Rf_*\C$是各个${}^pH^iRf_*\C[-i]$的直和，也不知道如何证明\term{perverse-sheaf}{反常层}${}^pH^iRf_*\C$\term{semisimple}{半单}，或对$Rf_*(\C[\dim X])$成立的6.2.10的类似命题。
